% Original analytic teaching diagram, 2026-11-12.
% Three-point parabola: X uniform on {-1,0,1}, Y = X^2.  Left: joint probability table P(X=x,Y=y);
% right: product of the marginals P(X=x)P(Y=y).  Cells with positive mass are light blue; zero cells of the
% joint table that are positive in the product are outlined in red.
\definecolor{VB}{HTML}{7188A3}\definecolor{VBF}{HTML}{E5EAF0}
\definecolor{VG}{HTML}{829C8B}\definecolor{VGF}{HTML}{E8EEE9}
\definecolor{VY}{HTML}{B4A16F}\definecolor{VYF}{HTML}{F1ECDD}
\definecolor{VR}{HTML}{A7787D}\definecolor{VRF}{HTML}{F0E5E6}
\definecolor{VP}{HTML}{9588A5}\definecolor{VPF}{HTML}{ECE8F0}
\definecolor{VI}{HTML}{3B4252}\definecolor{VM}{HTML}{6B7280}\definecolor{VL}{HTML}{4C4D4F}
\begin{tikzpicture}[x=1mm,y=1mm,>=Stealth,font=\figurefont\fontsize{8}{10}\selectfont,
lab/.style={text=VM,font=\figurefont\fontsize{7.5}{9}\selectfont},
mass/.style={draw=VB,fill=VBF,line width=.5pt},
zero/.style={draw=VL,fill=white,line width=.4pt},
miss/.style={draw=VR,line width=.9pt}]
% ---------- joint table ----------
\begin{scope}
\foreach \i/\x in {0/-1,1/0,2/1}{\node[lab] at (6+12*\i,25) {$x=\x$};}
\node[lab,anchor=east] at (-1,17) {$y=1$};
\node[lab,anchor=east] at (-1,7) {$y=0$};
% row y=1 : masses 1/3, 0, 1/3
\draw[mass] (0,12) rectangle (12,22); \node[text=VI] at (6,17) {$1/3$};
\draw[zero] (12,12) rectangle (24,22); \node[text=VI] at (18,17) {$0$};
\draw[mass] (24,12) rectangle (36,22); \node[text=VI] at (30,17) {$1/3$};
% row y=0 : 0, 1/3, 0
\draw[zero] (0,2) rectangle (12,12); \node[text=VI] at (6,7) {$0$};
\draw[mass] (12,2) rectangle (24,12); \node[text=VI] at (18,7) {$1/3$};
\draw[zero] (24,2) rectangle (36,12); \node[text=VI] at (30,7) {$0$};
\draw[miss] (12,12) rectangle (24,22);
\draw[miss] (0,2) rectangle (12,12);
\draw[miss] (24,2) rectangle (36,12);
\node[text=VI] at (18,-3) {$\mathbb{P}(X=x,\,Y=y)$};
\end{scope}
% ---------- product table ----------
\begin{scope}[xshift=56mm]
\foreach \i/\x in {0/-1,1/0,2/1}{\node[lab] at (6+12*\i,25) {$x=\x$};}
\node[lab,anchor=east] at (-1,17) {$y=1$};
\node[lab,anchor=east] at (-1,7) {$y=0$};
\foreach \i in {0,1,2}{
  \draw[mass] (12*\i,12) rectangle (12*\i+12,22); \node[text=VI] at (12*\i+6,17) {$2/9$};
  \draw[mass] (12*\i,2) rectangle (12*\i+12,12); \node[text=VI] at (12*\i+6,7) {$1/9$};
}
\node[text=VI] at (18,-3) {$\mathbb{P}(X=x)\,\mathbb{P}(Y=y)$};
\end{scope}
\end{tikzpicture}
